\section{Análisis teórico}

Sea \(T\) el texto y \(P\) el patrón. Denotamos sus longitudes por
\(n = |T|\) y \(m = |P|\), respectivamente. Excluimos el patrón vacío
(\(m = 0\)), cuyo número de coincidencias depende de la convención adoptada.
Analizamos el caso usual \(1 \leq m \leq n\); si \(m > n\), las implementaciones devuelven cero
inmediatamente. Por \textit{búsqueda} entendemos contar todas las ocurrencias,
incluidas las superpuestas. Usamos notación \(O(\cdot)\) para expresar cotas
superiores y \(\Theta(\cdot)\) cuando la cota es ajustada. El espacio se expresa
en palabras y no incluye la entrada. Como los métodos devuelven un contador, el
espacio de salida es \(O(1)\); si devolvieran todas las posiciones, habría que
sumar \(O(z)\), donde \(z\) es el número de ocurrencias.

\subsection{Búsqueda por fuerza bruta}

El algoritmo prueba cada posición de \(T\) y compara desde allí los caracteres
de \(P\), avanzando una posición tras cada intento.
\begin{itemize}
    \item \textbf{Preparación:} \(\Theta(1)\); no construye estructuras previas.
    \item \textbf{Mejor caso de búsqueda:} \(\Theta(n - m + 1)\), cuando cada
    intento falla al comparar el primer carácter. Una coincidencia inicial puede
    añadir hasta \(O(m)\) comparaciones, pero el algoritmo debe revisar las
    posiciones restantes para contar todas las ocurrencias.
    \item \textbf{Peor caso de búsqueda:} \(\Theta((n - m + 1)m)\), si cada ventana
    coincide en casi todo el patrón antes de fallar o confirmarse.
    \item \textbf{Memoria auxiliar:} \(\Theta(1)\).
\end{itemize}

\subsection{Knuth--Morris--Pratt (KMP)}

KMP construye una tabla que indica qué prefijo del patrón puede conservarse
después de una discrepancia. Así desplaza el patrón sin retroceder en el texto.
La implementación construye la tabla en cada llamada y luego
recorre el texto; tras una coincidencia, también usa la tabla para conservar
posibles coincidencias superpuestas.

\begin{itemize}
    \item \textbf{Preparación e inicialización:} \(\Theta(m)\) tiempo y
    memoria para construir la tabla de prefijos.
    \item \textbf{Mejor caso de búsqueda:} \(\Theta(n)\).
    \item \textbf{Peor caso de búsqueda:} \(\Theta(n)\); el recorrido avanza por
    el texto y los retrocesos del patrón son amortizadamente lineales.
    Incluida la preparación, el tiempo es \(\Theta(n+m)\).
    \item \textbf{Memoria auxiliar:} \(\Theta(m)\), por la tabla.
\end{itemize}

\subsection{Rabin--Karp}

Rabin--Karp calcula el hash del patrón \(P\) y el de la primera ventana de
\(T\) durante la preparación. La búsqueda verifica esa primera ventana y luego
desplaza la ventana, actualizando su hash en tiempo constante. Compara
caracteres para confirmar cada coincidencia de hash.

\begin{itemize}
    \item \textbf{Preparación e inicialización:} \(\Theta(m)\) tiempo para
    calcular ambos hashes: el de \(P\) y el de la primera ventana de \(T\).
    \item \textbf{Mejor caso de búsqueda:} \(\Theta(n - m + 1)\), cuando el
    hash de \(P\) difiere del hash de todas las ventanas examinadas. No hay
    verificaciones carácter a carácter y cada ventana se procesa en tiempo
    constante.
    \item \textbf{Peor caso de búsqueda:} \(\Theta((n - m + 1)m)\), si muchas
    ventanas requieren una verificación completa, por colisiones o coincidencias
    reales superpuestas.
    \item \textbf{Memoria auxiliar:} \(\Theta(1)\).
\end{itemize}

El peor caso puede ocurrir por colisiones de hash o por coincidencias reales en
muchas ventanas. Como se usan parámetros de hash fijos y no se supone una
distribución aleatoria de las entradas, no se afirma una garantía de tiempo
esperado. Se asume que las operaciones con el hash modular cuestan tiempo
constante.

\subsection{Suffix Array}

La estructura ordena los sufijos de \(T\) mediante duplicación de prefijos:
en cada ronda duplica el largo del prefijo ordenado y actualiza clases con
conteo. Para contar un patrón, realiza dos búsquedas binarias que delimitan
    los sufijos que comienzan con \(P\); no recorre individualmente las
\(z\) ocurrencias.

\begin{itemize}
    \item \textbf{Preparación (construcción):} \(O(n \log n)\) tiempo en el
    peor caso y \(\Theta(n)\) memoria, mediante duplicación de prefijos y
    ordenamiento por conteo en cada ronda. Esta implementación ejecuta todas
    las rondas, por lo que su construcción toma \(\Theta(n \log n)\).
    \item \textbf{Mejor caso de búsqueda:} \(\Theta(\log n)\), si las
    comparaciones binarias terminan al revisar el primer carácter.
    \item \textbf{Peor caso de búsqueda:} \(O(m \log n)\), pues cada comparación
    puede revisar hasta \(m\) caracteres. Es la búsqueda binaria básica, sin
    optimizaciones basadas en LCP; contar el intervalo cuesta \(O(1)\).
    \item \textbf{Memoria:} \(O(1)\) auxiliar por consulta; el índice
    persistente ocupa \(\Theta(n)\) y la memoria temporal de construcción,
    \(\Theta(n)\).
\end{itemize}

La memoria depende de los largos de entrada, no de su contenido; por ello no
cambia entre mejor y peor caso.

\subsection{Comparación de complejidades}

Los tiempos de búsqueda incluyen el conteo. En Rabin--Karp, la preparación
incluye el hash de \(P\) y el de la primera ventana de \(T\). La preparación del
Suffix Array corresponde a su construcción; se distingue el espacio auxiliar
por consulta del índice persistente y de la memoria temporal de construcción.

\begin{center}
\small
\begin{tabular}{@{}lcccc@{}}
\toprule
\textbf{Método} & \textbf{Preparación} & \textbf{Mejor búsqueda} &
\textbf{Peor búsqueda} & \shortstack{\textbf{Memoria auxiliar /}\\\textbf{índice}} \\
\midrule
Fuerza bruta & \(\Theta(1)\) & \(\Theta(n - m + 1)\) & \(\Theta((n - m + 1)m)\) & \(\Theta(1)\) \\
KMP & \(\Theta(m)\) & \(\Theta(n)\) & \(\Theta(n)\) & \(\Theta(m)\) \\
Rabin--Karp & \(\Theta(m)\) & \(\Theta(n - m + 1)\) & \(\Theta((n - m + 1)m)\) & \(\Theta(1)\) \\
Suffix Array & \(\Theta(n \log n)\) & \(\Theta(\log n)\) & \(O(m \log n)\) & \shortstack{\(O(1)\) consulta;\\\(\Theta(n)\) índice persistente;\\\(\Theta(n)\) temporal al construir} \\
\bottomrule
\end{tabular}
\end{center}
